Prompt tone is easy to change, common in real use, and poorly understood as a control variable in reasoning evaluations. Prior work studies politeness, emotional framing, and sycophancy, but it does not directly test whether skeptical or judgmental language creates a useful ``perform under scrutiny'' effect on objective tasks. We examine this question with a repeated-measures benchmark that varies only user tone. We ran 1,600 Chat Completions API calls on \gptfourone and \gptfouronemini across 40-item subsets of \gsm, \csqa, and \truthfulqa, under four prompt conditions and with an additional misleading-hint manipulation on the two multiple-choice datasets.

Our main result is a weak, model-dependent sweet spot. On pooled no-hint evaluations, mild skeptical questioning improves \gptfourone accuracy from 64.2\% to 69.2\%, a gain of 5.0 percentage points with a paired bootstrap 95\% confidence interval of [+0.8, +10.0]. Harsher judgment does not add further benefit. Under misleading wrong hints, the gain disappears or reverses, and direct agreement with the wrong hint remains low across all conditions. Across both models, judgmental phrasing increases rationale length much more reliably than it increases accuracy, while explicit apology markers never appear in any of the 1,600 responses.

These results suggest that skeptical prompting can sometimes sharpen answers, but harsh ``judging'' is not a reliable method for better reasoning. Its strongest and most reproducible effect is longer and more volatile visible reasoning rather than higher correctness.
